% ID: SEL-D0061-OA-E03
% ESTADO: validada. S061.
Sean $B$ «va en autobús» y $T$ «llega tarde». El coche tiene probabilidad $0{,}20$ y, si va en coche, llega tarde con probabilidad $0{,}90$.

\textbf{a)} $P(B\cap T^c)=0{,}80\cdot0{,}80=\boxed{0{,}64}$.

\textbf{b)} Por probabilidad total, $P(T)=0{,}80\cdot0{,}20+0{,}20\cdot0{,}90=\boxed{0{,}34}$.

\textbf{c)} El camino «va en coche y llega a tiempo» tiene probabilidad $0{,}20\cdot0{,}10=0{,}02$. Al dividir por $P(T^c)=0{,}66$, obtenemos $\boxed{P(B^c\mid T^c)=1/33}$.
